\documentclass{article}
\usepackage{/globals/de}
\usepackage{/globals/anki}
\ankiprefix{branchen}
\ankitopic{Branchen}
\ankishow
\begin{document}
\section{Branchen}
\subsection{Industrie}
Mithilfe von Lagererhaltungssystemen kann die Lagerung der zu verarbeitenden Materialien geplant werden.

Die Materialbewertung ist der Prozess, dass der Wert der Materialien berechnet wird.
\ankinote{materialbewertung}{Industrie Materialbewertung}{Wert der Materialien berechnen}

Bei der Bestelldisposition wird sichergestellt, dass immer genug Materialien vorhanden sind. Eine elektronische Lagerbestandsführung ist trivial, solange sich der Lagerstand nicht undokumentiert verändert.
\ankinote{bestelldisposition}{Industrie Bestelldisposition}{Sicherstellen dass genug Lagerbestand vorhanden ist}

Wird sich das Lager zeitlich dargestellt, sinkt der Lagerbestand konstant; die Materialien werden gebraucht. Nachdem eine Bestellgrenze erreicht wird, werden neue Materialien bestellt, welche nach der Beschaffungszeit am Liefertermin ankommen. Der Lagerbestand steigt um die Losgröße.
\ankinote{lager-zeitlich}{Lager zeitlich}{}

Die gesamte Inventur wird durch vollständige oder stichprobeartige Zählung bestimmt.

Modernisierte physische Produktionsverfahren in der Industrie können als Industrie 4.0 zusammengefasst werden.
\ankinote{modern}{Moderne Industrie}{Industrie 4.0}

\subsection{Dienstleistungen}
In der Dienstleistungebranche ist das verkaufte Produkt immateriell. Dies ist auch der größte Unterschied zur Industrie; es geht weniger um Materialflüsse und Lagerungen, sondern um Informationsflüsse. Dabei sind immer externe Faktoren beteiligt; \idr Kunden.

Arbeitsprozesse können am Ort unterschieden werden.
\begin{description}
 \item[Backoffice] Alle Prozesse die ohne direkten Kundenkontakt stattfinden. Dies kann durch einen Büroarbeiter oder durch eine Automatisierung stattfinden.
 \item[Frontoffice] Alle Prozesse mit direkten Kundenkontakt.

Frontofficeaufgaben passieren klassisch durch Arbeiter, können aber auch durch Self-Service Systeme automatisiert werden; eine DMZ-Platform auf die Kunden zugriff haben, die dessen Anfragen in das Unternehmensnetzwerk weiterleitet.
\end{description}
\ankinote{dienstleistungen}{Dienstleistungen}{Backoffice, Frontoffice}
\ankinote{backoffice}{Backoffice}{Kein Kundenkontakt; Büro oder Automatisierung}
\ankinote{frontoffice}{Frontoffice}{Kundenkontakt; Arbeiter oder Website}

\subsection{Finanzen}
IT-Systeme bieten Schnittstellen zu Banken und helfen bei bankeninternen Zahlungsprozessen. Sie können automatisiert ausgedruckte Belege einscannen und lesen.

Kredite können standardisiert bearbeitet werden.

Homogene Güter können über Handelsplattformen gehandelt werden.
\ankinote{finanzen}{Finanzen}{Schnittstellen, Standardisierung, Handel}

\subsection{Gesundheitswesen}
Akten können elektronisch gespeichert werden.

Es können vereinfacht Massen an Daten gesammelt und ausgewertet werden.
\ankinote{gesundheitswesen}{Gesundheitswesen}{Daten, Websiten, elektronische Akten}

\subsection{Medien}
Medien können digital verbreitet und verwaltet werden, mit Content Management Systemen.

Die digitalisierung ermöglicht auch erstmals das moderne Paradigma des Fokus auf User Generated Content.
\ankinote{medien}{Medien}{Verbreitung, Verwaltung (CMS), UGC}

\end{document}